\clearpage\pagestyle{otMain}\setlength{\otBlocEspacement}{2.5mm}\fontsize{11}{15}\selectfont\corrigesChapitre{1}{Multiplier et diviser des nombres relatifs}
\begin{blocLivre}{cor-1-1}\begin{corrige}{1}{Un produit positif}\demarche{Déterminer le signe avant de multiplier.}
Les deux facteurs ont le même signe, donc le résultat est positif. $7\times6=42$, ainsi $(-7)\times(-6)=\res{42}$.
\end{corrige}\end{blocLivre}
\begin{blocLivre}{cor-1-2}\begin{corrige}{2}{Un produit négatif}\demarche{Séparer signe et distance à zéro.}
Les signes sont différents : le produit est négatif. $2{,}4\times5=12$, donc le résultat est $\res{-12}$.
\end{corrige}\end{blocLivre}
\begin{blocLivre}{cor-1-3}\begin{corrige}{3}{Un quotient}\demarche{Appliquer la règle des signes à chaque quotient.}
Le premier quotient est positif : $54:9=\res6$. Le second est négatif : $18:4=4{,}5$, donc $\res{-4{,}5}$. Contrôle : $(-4{,}5)\times(-4)=18$.
\end{corrige}\end{blocLivre}
\begin{blocLivre}{cor-1-4}\begin{corrige}{4}{Opposé ou produit}\demarche{Identifier l’opération écrite entre les nombres.}
Soustraire $-3$ revient à ajouter 3 : $-8+3=\res{-5}$. Pour le produit, deux facteurs négatifs donnent un résultat positif : $\res{24}$.
\end{corrige}\end{blocLivre}
\begin{blocLivre}{cor-1-5}\begin{corrige}{5}{Nommer puis calculer}\demarche{Repérer l’opération effectuée en dernier.}
Dans A, la dernière opération est l’addition : A est une $\res{\text{somme}}$, $A=-3-28=\res{-31}$. B est un $\res{\text{produit}}$ : $B=-5\times(-5)=\res{25}$. C est un $\res{\text{quotient}}$ : $C=18:(-3)=\res{-6}$.
\end{corrige}\end{blocLivre}
\begin{blocLivre}{cor-1-6}\begin{corrige}{6}{Compter les signes}\demarche{Compter les facteurs négatifs non nuls.}
Il y a trois facteurs négatifs : le produit est négatif. Les distances à zéro donnent $2\times5\times3\times7=210$. Résultat : $\res{-210}$.
\end{corrige}\end{blocLivre}
\begin{blocLivre}{cor-1-7}\begin{corrige}{7}{Respecter les priorités}\demarche{Ne pas déplacer une parenthèse.}
$A=14-(-6)=\res{20}$. Dans $B$, les parenthèses passent d’abord : $B=11\times(-2)=\res{-22}$.
\end{corrige}\end{blocLivre}
\begin{blocLivre}{cor-1-8}\begin{corrige}{8}{Une expression imbriquée}\demarche{Commencer par la parenthèse la plus intérieure.}
$5-8=-3$, puis $-2\times(-3)=6$. Le quotient vaut donc $18:6=\res3$. Le crochet entier doit être calculé avant la division.
\end{corrige}\end{blocLivre}
\begin{blocLivre}{cor-1-9}\begin{corrige}{9}{Des multiplications à trou}\demarche{Retrouver un facteur en divisant.}
Le nombre cherché est un quotient. $42:(-6)=\res{-7}$ ; $(-10):(-2{,}5)=\res4$ ; $(-3):7=\res{-\frac37}$. Contrôle : $7\times\left(-\frac37\right)=-3$, par définition du quotient.
\end{corrige}\end{blocLivre}
\begin{blocLivre}{cor-1-10}\begin{corrige}{10}{Une erreur de calcul}\demarche{Conserver l’ordre de gauche à droite.}
Produits et quotients ont la même priorité et se calculent de gauche à droite : $24:(-6)=-4$, puis $-4\times2=\res{-8}$. Nora a ajouté des parenthèses absentes.
\end{corrige}\end{blocLivre}
\begin{blocLivre}{cor-1-11}\begin{corrige}{11}{Évolution de température}\demarche{Additionner la variation à la valeur initiale.}
La variation totale est $5\times(-1{,}6)=-8\,^\circ$C. On l’ajoute à la température initiale : $3-8=\res{-5\,^\circ\mathrm C}$. Une variation négative représente une baisse.
\end{corrige}\end{blocLivre}
\begin{blocLivre}{cor-1-12}\begin{corrige}{12}{Un programme de calcul}\demarche{Écrire chaque étape, puis les emboîter.}
On obtient successivement $(-4)\times(-3)=12$, puis $12+7=19$, puis $19\times(-2)=\res{-38}$. L’expression est $\res{[(-4)\times(-3)+7]\times(-2)}$ : la dernière opération est une multiplication, c’est un $\res{\text{produit}}$. Les crochets traduisent l’ordre des étapes.
\end{corrige}\end{blocLivre}
\begin{blocLivre}{cor-1-13}\begin{corrige}{13}{Un nombre inconnu}\demarche{Utiliser le quotient pour retrouver un facteur.}
On effectue l’opération inverse : $15:(-2{,}5)=\res{-6}$. Vérification : $(-2{,}5)\times(-6)=15$. Le signe négatif est nécessaire pour obtenir un produit positif.
\end{corrige}\end{blocLivre}
\begin{blocLivre}{cor-1-14}\begin{corrige}{14}{La machine à signes}\demarche{Traiter le cas nul séparément.}
19 est impair : le produit est négatif. Remplacer un facteur par son opposé change une seule fois le signe, donc le produit devient positif. Avec un facteur nul, tout le produit vaut $\res0$ : on ne lui applique plus la règle pair-impair.
\end{corrige}\end{blocLivre}
\begin{blocLivre}{cor-1-15}\begin{corrige}{15}{Un calcul à contrôler}\demarche{Calculer séparément les deux termes principaux.}
$-5+8=3$, puis $4-2\times3=-2$. Le premier terme vaut $(-3)\times(-2)=6$ et le second $-3$. Donc $C=6-3=\res3$. Le premier terme positif dépasse en valeur le second, négatif : le signe positif est cohérent.
\end{corrige}\end{blocLivre}
